\section{引言}
药物在不同溶剂中的溶解度影响溶剂筛选、结晶分离与制剂开发。实验难以覆盖全部药物、溶剂和温度组合，因此结构驱动的预测可为候选条件筛选提供参考\cite{bigsoldb2025,fastsolv2025}。实际数据中，同一药物常在多个条件下重复出现，逐行样本量也不能等同于独立分子的覆盖范围。

本文关注的任务是预测未见药物母体的条件溶解度。按记录随机划分可能让同一母体同时进入训练和验证，夸大对新药物的推广能力\cite{leakage2023}。另一方面，手工描述符、分子交互与温度信息之间存在非线性关系，仅增加模型结构复杂度未必能改善有限母体上的泛化。

TabPFN通过表格基础模型为小数据回归提供了不同于常规从零训练的方法\cite{tabpfn2025}；ExtraTrees以随机化树集成建模数值特征间的非线性关系\cite{extratrees2006}。两种模型的学习方式不同，因而值得在一致的分组协议下检验其互补性。本文以TabPFN＋ExtraTrees凸组合为主方法，通过内层母体验证同时组织基模型选择与融合选权。方法采用已有模型的任务化组合，不将异质融合本身宣称为原创算法。

本文首先分析数据覆盖与采样不均，随后给出双分支表示、母体加权树训练和折外选权流程，并与SOLAR结构--温度分解、通用回归器及多分支组合比较。模型比较同时考察母体与逐行误差，以检验融合收益及其在不同药物间的稳定性。
